% ======================================================================
%  Guía de laboratorio — DriveGuard Lab
%  Curso: Deep Learning y Agentes Inteligentes
%  Docente: Ing. Julián René Muñoz Burbano
%  Compilar con: pdflatex Guia_Laboratorio_DriveGuard.tex  (dos veces)
%  También compila en Overleaf (pdfLaTeX).
% ======================================================================
\documentclass[11pt,a4paper]{article}

\usepackage[utf8]{inputenc}
\usepackage[T1]{fontenc}
\IfFileExists{lmodern.sty}{\usepackage{lmodern}}{}
\usepackage[spanish,es-tabla,es-nodecimaldot]{babel}
\usepackage[margin=2.3cm,headheight=14pt]{geometry}
\usepackage{microtype}
\usepackage{amsmath,amssymb}
\usepackage{siunitx}
\sisetup{output-decimal-marker={,},group-minimum-digits=4}
\usepackage{graphicx}
\usepackage[table]{xcolor}
\usepackage{booktabs,tabularx,array,longtable}
\usepackage{enumitem}
\usepackage{fancyhdr}
\usepackage{float}
\usepackage{listings}
\usepackage[most]{tcolorbox}
\usepackage{tikz}
\usetikzlibrary{arrows.meta,positioning,shapes.geometric,calc,fit,backgrounds}
\usepackage[hidelinks]{hyperref}

% ---------------------------------------------------------------- colores
\definecolor{azul}{HTML}{1F4E79}
\definecolor{azulclaro}{HTML}{E8F0F8}
\definecolor{verde}{HTML}{2E9E4F}
\definecolor{amarillo}{HTML}{E6B800}
\definecolor{naranja}{HTML}{F07F13}
\definecolor{rojo}{HTML}{D62F2F}
\definecolor{gris}{HTML}{F4F5F7}
\definecolor{grisborde}{HTML}{C9CED6}

\hypersetup{colorlinks=true,linkcolor=azul,urlcolor=azul,citecolor=azul}

% ---------------------------------------------------------------- encabezado
\pagestyle{fancy}
\fancyhf{}
\fancyhead[L]{\small\color{azul} DriveGuard Lab — Guía de laboratorio}
\fancyhead[R]{\small\color{azul} Deep Learning y Agentes Inteligentes}
\fancyfoot[C]{\small \thepage}
\renewcommand{\headrulewidth}{0.4pt}

\setlength{\parskip}{4pt}
\setlength{\emergencystretch}{3em}
\setlist{itemsep=2pt,topsep=3pt}
\renewcommand{\arraystretch}{1.25}

% ---------------------------------------------------------------- código
\lstdefinestyle{consola}{
  basicstyle=\ttfamily\small, backgroundcolor=\color{gris}, frame=single,
  rulecolor=\color{grisborde}, framesep=4pt, xleftmargin=6pt, xrightmargin=6pt,
  breaklines=true, columns=fullflexible, keepspaces=true, showstringspaces=false,
  literate={á}{{\'a}}1 {é}{{\'e}}1 {í}{{\'i}}1 {ó}{{\'o}}1 {ú}{{\'u}}1 {ñ}{{\~n}}1
           {Á}{{\'A}}1 {É}{{\'E}}1 {Í}{{\'I}}1 {Ó}{{\'O}}1 {Ú}{{\'U}}1 {Ñ}{{\~N}}1 {°}{{$^\circ$}}1
}
\lstdefinestyle{python}{style=consola, language=Python,
  keywordstyle=\color{azul}\bfseries, commentstyle=\color{gray}\itshape, stringstyle=\color{rojo!80!black}}
\lstset{style=consola}

% ---------------------------------------------------------------- cajas
\newtcolorbox{objetivo}[1][]{colback=azulclaro,colframe=azul,fonttitle=\bfseries,title=#1,arc=2mm,boxrule=0.6pt}
\newtcolorbox{aviso}[1][Atención]{colback=rojo!6,colframe=rojo,fonttitle=\bfseries,title=#1,arc=2mm,boxrule=0.6pt}
\newtcolorbox{nota}[1][Nota]{colback=amarillo!10,colframe=amarillo!80!black,fonttitle=\bfseries,title=#1,arc=2mm,boxrule=0.6pt}
\newtcolorbox{verificacion}[1][Punto de verificación]{colback=verde!7,colframe=verde,fonttitle=\bfseries,title=#1,arc=2mm,boxrule=0.6pt}

\newcommand{\codigoguion}{\textunderscore\allowbreak}
\newcommand{\codigo}[1]{{\small\ttfamily\let\_\codigoguion #1}}
\newcommand{\nivel}[2]{\tikz[baseline=-0.6ex]\node[fill=#1,text=white,rounded corners=2pt,inner sep=2pt,font=\scriptsize\bfseries]{#2};}

\begin{document}

% ======================================================================
%  PORTADA
% ======================================================================
\begin{titlepage}
\begin{tikzpicture}[remember picture,overlay]
  \fill[azul] (current page.north west) rectangle ([yshift=-7.2cm]current page.north east);
  \node[anchor=north west,text=white,font=\Large\bfseries] at ([xshift=2.2cm,yshift=-1.8cm]current page.north west)
    {Deep Learning y Agentes Inteligentes};
  \node[anchor=north west,text=white,font=\large] at ([xshift=2.2cm,yshift=-2.6cm]current page.north west)
    {Guía de laboratorio};
  \node[anchor=north west,text=white,font=\fontsize{34}{40}\selectfont\bfseries] at ([xshift=2.2cm,yshift=-3.8cm]current page.north west)
    {DriveGuard Lab};
  \node[anchor=north west,text=white!85!azul,font=\Large] at ([xshift=2.2cm,yshift=-5.3cm]current page.north west)
    {Detección de somnolencia y riesgo del conductor};
\end{tikzpicture}

\vspace*{7.2cm}

\begin{center}
\begin{tikzpicture}[scale=0.95]
  % semáforo de niveles
  \foreach \x/\c/\t in {0/verde/NORMAL,3.2/amarillo/BAJO,6.4/naranja/MEDIO,9.6/rojo/ALTO}{
    \fill[\c,rounded corners=4pt] (\x,0) rectangle +(2.8,1.1);
    \node[text=white,font=\bfseries] at (\x+1.4,0.55) {\t};
  }
  \draw[-{Latex[length=3mm]},thick,azul] (0,-0.5) -- (12.4,-0.5) node[midway,below,font=\small] {nivel de riesgo};
\end{tikzpicture}
\end{center}

\vspace{1.2cm}
\begin{center}
\begin{tabular}{>{\bfseries}r l}
Curso: & Deep Learning y Agentes Inteligentes\\
Docente: & Ing. Julián René Muñoz Burbano\\
Proyecto: & DriveGuard Lab v1.0 (continuación de SentimentLab V3)\\
Duración: & 4 sesiones de 2 horas\\
Modalidad: & Individual o en parejas\\
Entorno: & Anaconda · Jupyter · Streamlit · FastAPI · MongoDB\\
\end{tabular}
\end{center}

\vfill
\begin{aviso}[Aviso de seguridad]
Este sistema es un \textbf{prototipo académico}. No sustituye sistemas de seguridad vehicular certificados.
Todas las pruebas se realizan frente al computador, con el vehículo detenido o con videos grabados.
\textbf{Nunca} lo pruebe mientras conduce en vía pública.
\end{aviso}
\end{titlepage}

\tableofcontents
\newpage

% ======================================================================
\section{Introducción}
% ======================================================================
La somnolencia al volante es una de las causas de siniestros viales más difíciles de detectar: el conductor
suele no darse cuenta de que está a punto de dormirse. Un \emph{microsueño} de dos segundos a
\SI{80}{km/h} equivale a recorrer más de \SI{44}{m} sin control del vehículo.

En este laboratorio se construye \textbf{DriveGuard Lab}, un sistema que observa al conductor con una
cámara web y estima su nivel de riesgo a partir de señales visuales: cuánto tiempo mantiene los ojos
cerrados, con qué frecuencia bosteza, si cabecea o gira la cabeza y qué expresión facial aparenta. Cuando
detecta riesgo, hace sonar una alarma; si el conductor no responde, avisa automáticamente a sus contactos
de emergencia por correo electrónico y Telegram.

El proyecto continúa \emph{SentimentLab V3}: conserva su arquitectura (\codigo{core/}, API con FastAPI,
interfaz con Streamlit, MongoDB, notebooks y pruebas) y reutiliza su modelo de expresiones faciales.

% ======================================================================
\section{Objetivos}
% ======================================================================
\begin{objetivo}[Objetivo general]
Diseñar, implementar y evaluar un sistema de visión por computador que detecte señales de somnolencia y
distracción en un conductor, genere alertas sonoras locales y ejecute un protocolo de emergencia,
integrando modelos preentrenados, un motor de reglas, una API REST y una base de datos NoSQL.
\end{objetivo}

\textbf{Objetivos específicos.} Al finalizar el laboratorio, el estudiante será capaz de:
\begin{enumerate}
  \item Calcular e interpretar las métricas EAR, MAR y PERCLOS a partir de puntos faciales de MediaPipe.
  \item Explicar por qué un detector de rostros (Haar Cascade) no basta para medir somnolencia.
  \item Implementar un motor de riesgo con calibración por conductor, ventanas temporales e histéresis.
  \item Justificar el uso de hilos separados para video, clasificación de expresiones y alarma.
  \item Configurar y probar un protocolo de emergencia con notificaciones reales o simuladas.
  \item Entrenar una CNN para clasificar ojos abiertos y cerrados, y compararla con la regla del EAR.
  \item Evaluar el sistema con métricas (precisión, recall, F1, falsas alarmas, latencia) y analizar sus
        limitaciones técnicas y éticas.
\end{enumerate}

% ======================================================================
\section{Marco teórico}
% ======================================================================
\subsection{Puntos faciales con MediaPipe FaceLandmarker}
\emph{FaceLandmarker} es un modelo preentrenado de Google que, a partir de una imagen RGB, estima
\textbf{478 puntos tridimensionales} del rostro, 52 coeficientes de expresión (\emph{blendshapes}) y una
matriz de transformación $4\times4$ que describe la orientación de la cabeza. Funciona en tiempo real en
CPU (entre 10 y \SI{25}{ms} por fotograma en un portátil).

\begin{nota}[¿Por qué no basta el Haar Cascade de SentimentLab V3?]
El detector Haar Cascade devuelve solo el \emph{rectángulo} donde está la cara. Para saber si los ojos están
cerrados o la boca abierta se necesita la \emph{forma} de los párpados y los labios, es decir, puntos
faciales. Por eso DriveGuard Lab usa FaceLandmarker (API \emph{Tasks} de MediaPipe).
\end{nota}

\subsection{Eye Aspect Ratio (EAR)}
Soukupová y Čech \cite{soukupova2016} proponen medir la apertura del ojo con seis puntos: dos comisuras
($p_1$, $p_4$) y dos pares de puntos en los párpados superior ($p_2$, $p_3$) e inferior ($p_6$, $p_5$):
\begin{equation}
  \mathrm{EAR} = \frac{\lVert p_2 - p_6\rVert + \lVert p_3 - p_5\rVert}{2\,\lVert p_1 - p_4\rVert}
  \label{eq:ear}
\end{equation}

\begin{figure}[H]
\centering
\begin{tikzpicture}[scale=0.9,
  pt/.style={circle,fill=azul,inner sep=1.6pt},
  lbl/.style={font=\small}]
  % ojo abierto
  \begin{scope}
    \draw[thick,fill=azulclaro] (0,0) .. controls (1.2,1.1) and (2.8,1.1) .. (4,0) .. controls (2.8,-1.1) and (1.2,-1.1) .. (0,0);
    \fill[azul!70] (2,0) circle (0.55);
    \node[pt,label=left:{\small $p_1$}] (a) at (0,0) {};
    \node[pt,label=above:{\small $p_2$}] (b) at (1.3,0.8) {};
    \node[pt,label=above:{\small $p_3$}] (c) at (2.7,0.8) {};
    \node[pt,label=right:{\small $p_4$}] (d) at (4,0) {};
    \node[pt,label=below:{\small $p_5$}] (e) at (2.7,-0.8) {};
    \node[pt,label=below:{\small $p_6$}] (f) at (1.3,-0.8) {};
    \draw[rojo,thick,dashed] (b)--(f) (c)--(e);
    \draw[verde,thick,dashed] (a)--(d);
    \node[lbl] at (2,-1.9) {Ojo abierto: $\mathrm{EAR}\approx 0{,}30$};
  \end{scope}
  % ojo cerrado
  \begin{scope}[xshift=6.5cm]
    \draw[thick,fill=azulclaro] (0,0) .. controls (1.2,0.18) and (2.8,0.18) .. (4,0) .. controls (2.8,-0.18) and (1.2,-0.18) .. (0,0);
    \node[pt] (a) at (0,0) {}; \node[pt] (d) at (4,0) {};
    \node[pt] (b) at (1.3,0.13) {}; \node[pt] (c) at (2.7,0.13) {};
    \node[pt] (e) at (2.7,-0.13) {}; \node[pt] (f) at (1.3,-0.13) {};
    \draw[verde,thick,dashed] (a)--(d);
    \node[lbl] at (2,-1.9) {Ojo cerrado: $\mathrm{EAR}\approx 0{,}05$};
  \end{scope}
\end{tikzpicture}
\caption{Puntos del EAR. Las distancias verticales (rojo) se dividen por la horizontal (verde).}
\label{fig:ear}
\end{figure}

El EAR es \textbf{invariante a la escala}: no cambia si el conductor se acerca o se aleja de la cámara.
Un ojo abierto da valores entre 0,25 y 0,35; uno cerrado, menos de 0,15. En la malla de MediaPipe los
índices son \codigo{[33, 160, 158, 133, 153, 144]} (ojo derecho) y \codigo{[362, 385, 387, 263, 373, 380]}
(ojo izquierdo).

\subsection{Mouth Aspect Ratio (MAR) y bostezos}
De forma análoga, la apertura de la boca se mide con las comisuras internas ($c_1=78$, $c_2=308$) y tres
pares verticales $(81,178)$, $(13,14)$ y $(311,402)$:
\begin{equation}
  \mathrm{MAR} = \frac{\tfrac{1}{3}\sum_{k=1}^{3}\lVert v_k^{\text{sup}} - v_k^{\text{inf}}\rVert}{\lVert c_1 - c_2\rVert}
\end{equation}
Con la boca cerrada $\mathrm{MAR}\approx 0$–$0{,}1$; al hablar, $0{,}2$–$0{,}4$. Se cuenta un
\textbf{bostezo} cuando $\mathrm{MAR}>0{,}6$ durante al menos \SI{2}{s}.

\subsection{PERCLOS}
PERCLOS (\emph{PERcentage of eyelid CLOSure}) es la proporción de tiempo que los ojos permanecen cerrados en
una ventana deslizante \cite{wierwille1994}. Es uno de los indicadores de somnolencia más usados:
\begin{equation}
  \mathrm{PERCLOS} = \frac{\text{fotogramas con ojos cerrados en los últimos } 60\ \mathrm{s}}{\text{fotogramas con rostro en los últimos } 60\ \mathrm{s}}
\end{equation}
Una persona alerta parpadea unas 15–20 veces por minuto (PERCLOS $<10$\,\%). Valores superiores al 30\,\%
indican cansancio importante.

\subsection{Orientación de la cabeza}
De la matriz de transformación $R$ se obtienen tres ángulos: \emph{pitch} (inclinación arriba–abajo),
\emph{yaw} (giro izquierda–derecha) y \emph{roll}. Durante la calibración se toma la postura natural del
conductor como referencia; un \textbf{cabeceo} es una inclinación de más de $15^\circ$ respecto a esa
referencia y una \textbf{distracción}, un giro de más de $30^\circ$.

\subsection{Motor de riesgo}
El motor combina los indicadores en cuatro niveles (tabla~\ref{tab:niveles}).

\begin{table}[H]
\centering
\caption{Reglas por defecto del motor de riesgo (configurables en la aplicación).}
\label{tab:niveles}
\begin{tabularx}{\textwidth}{l X l}
\toprule
\textbf{Indicador} & \textbf{Regla} & \textbf{Nivel}\\
\midrule
Microsueño & ojos cerrados de forma continua $\geq$ \SI{1,5}{s} & \nivel{rojo}{ALTO}\\
Cabeza caída & inclinación $>15^\circ$ durante $\geq$ \SI{3}{s} & \nivel{rojo}{ALTO}\\
Cansancio (PERCLOS) & $>30$\,\% del último minuto con ojos cerrados & \nivel{naranja}{MEDIO}\\
Bostezos & $\geq 3$ bostezos en 5 minutos & \nivel{naranja}{MEDIO}\\
Cabeceos & $\geq 3$ cabeceos en 60 s & \nivel{naranja}{MEDIO}\\
Distracción & cabeza girada $>30^\circ$ o rostro ausente durante $>$ \SI{3}{s} & \nivel{naranja}{MEDIO}\\
Expresión negativa & $>70$\,\% de clasificaciones negativas en 30 s & \nivel{amarillo}{BAJO}\\
\bottomrule
\end{tabularx}
\end{table}

Tres mecanismos hacen que el motor sea robusto:
\begin{itemize}
  \item \textbf{Calibración por conductor.} Durante los primeros \SI{10}{s} se mide la mediana del EAR del
        conductor con los ojos abiertos (EAR base). El umbral de ojo cerrado es el 75\,\% de ese valor. Así
        se adapta a personas con ojos más grandes o más pequeños.
  \item \textbf{Suavizado.} El EAR se promedia en los últimos 3 fotogramas para reducir el ruido.
  \item \textbf{Histéresis.} El nivel sube de inmediato, pero solo baja tras \SI{2}{s} sin la condición. Esto
        evita que la alarma se encienda y apague con cada fotograma.
\end{itemize}

\subsection{Expresiones faciales: ViT y CNN}
Se reutiliza el clasificador \codigo{trpakov/vit-face-expression} de SentimentLab V3, un \emph{Vision
Transformer} (ViT) \cite{dosovitskiy2021} que divide la imagen en parches y los procesa con mecanismos de
atención. Una \textbf{CNN}, en cambio, aplica filtros convolucionales que extraen patrones locales. En el
laboratorio 04 se entrena una CNN pequeña para clasificar recortes de ojos.

\begin{nota}[Expresiones no son emociones]
Las etiquetas del ViT (enojo, tristeza, miedo\dots) describen \emph{apariencias} faciales, no emociones
internas. El modelo tiene sesgos de su conjunto de entrenamiento. Por eso, en DriveGuard Lab, la expresión
negativa solo genera un riesgo \emph{bajo}.
\end{nota}

\subsection{Arquitectura multihilo y alarma local}
\codigo{streamlit-webrtc} procesa cada fotograma en un \textbf{hilo de video} independiente del hilo de
Streamlit. Desde ese hilo \textbf{no se puede} usar \codigo{st.audio} ni modificar la página. Como la
aplicación se ejecuta en el mismo portátil que tiene la cámara, la alarma suena localmente con
\codigo{pygame} (o \codigo{winsound} en Windows) desde un hilo propio. El ViT, que tarda entre 150 y
\SI{600}{ms} en CPU, también corre en un hilo aparte y solo cada \codigo{EXPRESION\_INTERVALO\_S} segundos,
para no frenar la detección de somnolencia.

\subsection{Protocolo de emergencia}
Una cámara que mira al conductor \textbf{no puede detectar un choque}. En un vehículo real eso lo hace un
acelerómetro (sistemas como eCall). DriveGuard Lab define tres disparadores (figura~\ref{fig:estados}):
\begin{enumerate}
  \item \textbf{Conductor que no responde:} nivel ALTO y \SI{10}{s} más con ojos cerrados, cabeza caída o
        sin rostro visible.
  \item \textbf{Botón SOS:} envío inmediato.
  \item \textbf{Choque simulado:} botón o archivo CSV de acelerómetro con un pico mayor que \SI{4}{g}.
\end{enumerate}
En los casos 1 y 3 se muestra una cuenta regresiva de \SI{15}{s} con el botón \emph{Estoy bien} para
cancelar falsas alarmas.

\begin{figure}[H]
\centering
\begin{tikzpicture}[node distance=1.3cm and 2.2cm,
  est/.style={rectangle,rounded corners=3pt,draw=azul,fill=azulclaro,thick,minimum width=2.6cm,minimum height=0.9cm,align=center,font=\small},
  fin/.style={est,draw=rojo,fill=rojo!8},
  flecha/.style={-{Latex[length=2.5mm]},thick},
  lab/.style={font=\scriptsize,align=center,fill=white,inner sep=1pt}]
  \node[est] (mon) {Monitoreo};
  \node[est,right=3cm of mon] (alto) {Riesgo\\ALTO};
  \node[est,below=1.6cm of alto] (cuenta) {Cuenta regresiva\\15 s};
  \node[fin,right=2.6cm of cuenta] (envio) {Envío de alertas\\correo, Telegram};
  \draw[flecha] (mon) to[bend left=15] node[lab,above]{ojos cerrados $\geq$ 1,5 s} (alto);
  \draw[flecha] (alto) to[bend left=15] node[lab,below]{abre los ojos} (mon);
  \draw[flecha] (alto) -- node[lab,right]{10 s sin\\respuesta} (cuenta);
  \draw[flecha] (mon) |- node[lab,pos=0.25,right]{choque\\simulado} (cuenta);
  \draw[flecha] (cuenta) -- node[lab,above]{15 s} (envio);
  \draw[flecha] (cuenta.south) -- ++(0,-0.6) -| node[lab,pos=0.25,below]{``Estoy bien''} ([xshift=-1.2cm]mon.west) -- (mon.west);
  \draw[flecha] (mon.north) -- ++(0,0.9) -| node[lab,pos=0.25,above]{botón SOS (inmediato)} (envio.north);
\end{tikzpicture}
\caption{Estados del protocolo de emergencia.}
\label{fig:estados}
\end{figure}

% ======================================================================
\section{Arquitectura del sistema}
% ======================================================================
\begin{figure}[H]
\centering
\begin{tikzpicture}[scale=0.94,transform shape,
  caja/.style={rectangle,rounded corners=3pt,draw=azul,fill=white,thick,minimum height=0.85cm,align=center,font=\small,text width=3.1cm},
  flecha/.style={-{Latex[length=2.3mm]},thick,azul},
  grupo/.style={rectangle,rounded corners=5pt,draw=grisborde,fill=gris,inner sep=7pt}]
  \node[caja,fill=azulclaro] (cam) {Cámara web};
  \node[caja,right=0.9cm of cam] (mp) {MediaPipe\\EAR, MAR, pose};
  \node[caja,right=0.9cm of mp] (motor) {Motor de riesgo};
  \node[caja,below=0.7cm of mp] (vit) {ViT expresiones\\(hilo aparte)};
  \node[caja,below=0.7cm of motor] (alarma) {Alarma local\\pygame / winsound};
  \node[caja,above=0.8cm of motor] (panel) {Panel Streamlit\\semáforo, SOS};
  \node[caja,right=1.3cm of panel] (api) {API FastAPI};
  \node[caja,below=0.8cm of api,fill=verde!10] (db) {MongoDB};
  \node[caja,below=0.7cm of db,fill=rojo!8] (notif) {Correo, Telegram\\SMS (opcional)};
  \draw[flecha] (cam) -- (mp);
  \draw[flecha] (mp) -- (motor);
  \draw[flecha,dashed] (mp) -- (vit);
  \draw[flecha,dashed] (vit) -- (motor);
  \draw[flecha] (motor) -- (alarma);
  \draw[flecha] (motor) -- (panel);
  \draw[flecha] (panel) -- node[above,font=\scriptsize]{eventos} (api);
  \draw[flecha] (api) -- (db);
  \draw[flecha] (api.east) -- ++(0.4,0) |- (notif.east);
  \begin{scope}[on background layer]
    \node[grupo,fit=(mp)(motor)(vit)(alarma),label={[font=\scriptsize\bfseries,text=azul]below:core/ (hilo de video)}] {};
  \end{scope}
\end{tikzpicture}
\caption{Componentes principales de DriveGuard Lab.}
\end{figure}

\begin{table}[H]
\centering
\caption{Cambios respecto a SentimentLab V3 y su justificación.}
\begin{tabularx}{\textwidth}{>{\raggedright}p{3.6cm} X}
\toprule
\textbf{Cambio} & \textbf{Justificación}\\
\midrule
MediaPipe FaceLandmarker & El Haar Cascade solo encuentra la cara; para ojos y boca se necesitan puntos faciales.\\
Alarma local con pygame/winsound & El hilo de video de streamlit-webrtc no puede reproducir sonido en la página.\\
Protocolo de emergencia & Una cámara no detecta choques: se usan conductor que no responde, SOS y choque simulado.\\
Creación automática de \codigo{data/} & En el laboratorio 02 la carpeta no venía en el zip y la celda de cámara fallaba.\\
ViT en hilo aparte cada $N$ s & El ViT es lento en CPU; así no afecta la detección de somnolencia.\\
\bottomrule
\end{tabularx}
\end{table}

% ======================================================================
\section{Materiales y entorno}
% ======================================================================
\begin{itemize}
  \item Portátil con Windows 10/11 (o macOS/Linux), 8 GB de RAM como mínimo y cámara web.
  \item Anaconda o Miniconda con Python 3.11.
  \item MongoDB Community Server (servicio local) y, opcionalmente, MongoDB Compass.
  \item Navegador Chrome o Edge actualizado.
  \item Una cuenta de Gmail con verificación en dos pasos (para el correo) y Telegram (para los mensajes).
  \item Proyecto \codigo{DriveGuardLab.zip} entregado por el docente.
\end{itemize}

% ======================================================================
\section{Instalación}
% ======================================================================
Abra \textbf{Anaconda Prompt}, descomprima el proyecto y ubíquese en su carpeta.

\begin{lstlisting}
conda create -n DriveGuardLab python=3.11 -y
conda activate DriveGuardLab
python -m pip install -r requirements.txt
python -m ipykernel install --user --name DriveGuardLab --display-name "Python (DriveGuard Lab)"
copy .env.example .env
python scripts/descargar_modelos.py
python scripts/generar_recursos.py
python -m pytest
\end{lstlisting}

\begin{verificacion}
\codigo{pytest} debe terminar con \textbf{43 passed}. Si falla alguna importación, revise que el entorno
\codigo{DriveGuardLab} esté activado.
\end{verificacion}

Para ejecutar la aplicación se necesitan \textbf{dos consolas} con el entorno activado:
\begin{lstlisting}
:: Consola 1 - API
python -m uvicorn backend.main:app --reload --port 8000

:: Consola 2 - Interfaz
python -m streamlit run frontend/app.py
\end{lstlisting}

% ======================================================================
\section{Procedimiento}
% ======================================================================

% ----------------------------------------------------------------------
\subsection{Práctica 1 — Verificación del entorno (sesión 1, 40 min)}
\begin{enumerate}
  \item Abra \codigo{jupyter lab} desde la raíz del proyecto y ejecute
        \codigo{notebooks/00\_Verificacion\_Entorno.ipynb} con el kernel \emph{Python (DriveGuard Lab)}.
  \item Compruebe que todas las librerías aparecen con \codigo{[OK]} y que el modelo
        \codigo{face\_landmarker.task} se descargó.
  \item Cambie \codigo{SONAR = True} y escuche los dos sonidos de alarma.
  \item Cambie \codigo{USE\_CAMERA = True} y tome una captura. Observe que se guarda en
        \codigo{data/capturas/} sin errores: la carpeta se crea automáticamente.
  \item Verifique que MongoDB responde. Si no, inicie el servicio en \codigo{services.msc}.
\end{enumerate}

\begin{verificacion}
Captura guardada, sonidos audibles y mensaje \codigo{[OK] MongoDB responde}.
\end{verificacion}

% ----------------------------------------------------------------------
\subsection{Práctica 2 — EAR, MAR, PERCLOS y motor de riesgo (sesiones 1 y 2)}
Trabaje en \codigo{notebooks/03\_Somnolencia\_EAR\_MAR.ipynb}.
\begin{enumerate}
  \item \textbf{Secciones 2–4.} Ejecute las celdas de EAR y MAR. Verifique a mano, con la
        ecuación~\eqref{eq:ear}, el EAR del ojo abierto de la figura del notebook.
  \item \textbf{Sección 5.} Ejecute el motor con los siete escenarios sintéticos. Registre en la
        tabla~\ref{tab:escenarios} el nivel máximo y el instante en que se activa cada alerta.
  \item \textbf{Sección 6.} Ejecute el experimento de umbrales y complete la tabla~\ref{tab:umbrales}.
        Elija la configuración que considere más adecuada y justifíquela.
  \item \textbf{Sección 7.} Tome una foto suya y visualice los 478 puntos. Anote su EAR con los ojos
        abiertos, entrecerrados y cerrados.
  \item \textbf{Sección 8.} Grabe un video de 60–90 s con este guion:
        \begin{enumerate}[label=\alph*)]
          \item 0–10 s: mirar al frente (calibración).
          \item 10–25 s: parpadeo normal.
          \item 25–30 s: ojos cerrados \SI{3}{s}.
          \item 35–45 s: un bostezo largo.
          \item 50–60 s: tres cabeceos.
          \item 65–75 s: girar la cabeza a un lado \SI{5}{s}.
        \end{enumerate}
        Analícelo y compare las alertas detectadas con lo que realmente hizo.
\end{enumerate}

\begin{table}[H]
\centering
\caption{Registro de escenarios sintéticos (complete).}
\label{tab:escenarios}
\begin{tabularx}{\textwidth}{l X X X X}
\toprule
\textbf{Escenario} & \textbf{Nivel máximo} & \textbf{Instante (s)} & \textbf{Motivo} & \textbf{¿Emergencia?}\\
\midrule
normal & & & & \\
cansancio & & & & \\
microsueño & & & & \\
sin\_respuesta & & & & \\
distracción & & & & \\
cabeceos & & & & \\
rostro\_ausente & & & & \\
\bottomrule
\end{tabularx}
\end{table}

\begin{table}[H]
\centering
\caption{Experimento de umbrales: detección / falsas alarmas (\%).}
\label{tab:umbrales}
\begin{tabularx}{\textwidth}{c *{3}{>{\centering\arraybackslash}X}}
\toprule
\textbf{microsueno\_s} & \textbf{factor 0,65} & \textbf{factor 0,75} & \textbf{factor 0,85}\\
\midrule
0,5 s & \hspace{1.5cm}/ & \hspace{1.5cm}/ & \hspace{1.5cm}/ \\
1,0 s & \hspace{1.5cm}/ & \hspace{1.5cm}/ & \hspace{1.5cm}/ \\
1,5 s & \hspace{1.5cm}/ & \hspace{1.5cm}/ & \hspace{1.5cm}/ \\
2,0 s & \hspace{1.5cm}/ & \hspace{1.5cm}/ & \hspace{1.5cm}/ \\
\bottomrule
\end{tabularx}
\end{table}

% ----------------------------------------------------------------------
\subsection{Práctica 3 — Aplicación en vivo y protocolo de emergencia (sesión 3)}
\begin{enumerate}
  \item Inicie la API y la interfaz. En la barra lateral deben aparecer \emph{API conectada} y
        \emph{MongoDB conectado}.
  \item \textbf{Conductor.} Registre su nombre y al menos un contacto de emergencia (un compañero) con
        correo y \codigo{chat\_id} de Telegram.
  \item \textbf{Notificaciones.} Configure Gmail y Telegram en \codigo{.env} (ver README, sección 7).
        Reinicie la API y use \emph{Configuración $\rightarrow$ Enviar mensaje de prueba}. Si no tiene
        credenciales, active \codigo{NOTIFICACIONES\_SIMULADAS=true}.
  \item \textbf{Monitoreo en vivo.} Pulse \emph{START}, mire al frente durante la calibración y observe el
        semáforo. Pruebe, en este orden:
        \begin{enumerate}[label=\alph*)]
          \item parpadeos normales (debe permanecer en NORMAL);
          \item ojos cerrados \SI{2}{s} (ALTO y sirena);
          \item \emph{Estoy bien} (se silencia la alarma);
          \item tres bostezos largos (MEDIO);
          \item girar la cabeza \SI{4}{s} (distracción);
          \item ojos cerrados más de \SI{12}{s} (conductor que no responde y cuenta regresiva). Deje que
                termine la cuenta regresiva y confirme que el compañero recibe el mensaje.
        \end{enumerate}
  \item \textbf{SOS y choque.} Pulse \emph{SOS}. Luego abra \emph{Simulación de impacto}, cargue el archivo
        \codigo{acelerometro\_choque.csv} de la carpeta \codigo{data/ejemplos} y cancele la cuenta regresiva con
        \emph{Estoy bien}. Repita con \codigo{acelerometro\_sin\_choque.csv}.
  \item \textbf{Historial.} Revise la sesión: eventos por tipo, gráficas de EAR y PERCLOS y alertas enviadas.
        Exporte los eventos a CSV.
  \item \textbf{Rendimiento.} Anote los FPS con las expresiones activadas y desactivadas, y con
        \codigo{EXPRESION\_INTERVALO\_S} = 1, 2 y 5 s.
\end{enumerate}

\begin{verificacion}
El contacto recibe el mensaje con nombre, hora, motivo, nivel y ubicación. El historial muestra la sesión
con sus eventos y las alertas enviadas. Una segunda alerta en menos de 2 minutos aparece como
\emph{omitida\_cooldown}.
\end{verificacion}

\begin{table}[H]
\centering
\caption{Registro de pruebas en vivo (complete).}
\begin{tabularx}{\textwidth}{X c c c X}
\toprule
\textbf{Prueba} & \textbf{¿Detectó?} & \textbf{Nivel} & \textbf{Latencia (s)} & \textbf{Observaciones}\\
\midrule
Parpadeo normal & & & & \\
Ojos cerrados 2 s & & & & \\
Bostezos & & & & \\
Giro de cabeza & & & & \\
Sin respuesta 12 s & & & & \\
SOS & & & & \\
Choque simulado & & & & \\
\bottomrule
\end{tabularx}
\end{table}

% ----------------------------------------------------------------------
\subsection{Práctica 4 — CNN de ojos abiertos y cerrados (sesión 4)}
Trabaje en \codigo{notebooks/04\_CNN\_Ojos\_Abiertos\_Cerrados.ipynb}.
\begin{enumerate}
  \item Descargue el MRL Eye Dataset \cite{fusek2018} y descomprímalo en \codigo{data/ojos/mrl/}, o capture
        sus propios ojos con la sección 2 del notebook (\codigo{USE\_CAMERA = True}).
  \item Entrene la CNN (sección 5) y registre la curva de pérdida y la exactitud de validación.
  \item Evalúe en el conjunto de prueba: exactitud, F1 por clase y matriz de confusión.
  \item Mida el tiempo de inferencia por fotograma.
  \item Grabe un video con intervalos conocidos de ojos cerrados y compare EAR vs.\ CNN (sección 8).
  \item Active la CNN en la aplicación con \codigo{OJOS\_METODO=cnn} y repita la prueba de microsueño.
\end{enumerate}

\begin{aviso}[Modo demostración]
Si no hay datos, el notebook genera ojos sintéticos solo para comprobar que el código funciona. Esas
métricas no son válidas para el informe ni para usar la CNN en la aplicación.
\end{aviso}

% ======================================================================
\section{Preguntas de análisis}
% ======================================================================
\begin{enumerate}
  \item ¿Por qué el EAR es invariante a la distancia entre el conductor y la cámara? Demuéstrelo a partir de
        la ecuación~\eqref{eq:ear}.
  \item ¿Qué ocurre con el EAR de una persona con gafas oscuras? ¿Y de noche sin iluminación infrarroja?
  \item Explique el compromiso entre detección y falsas alarmas que observó en el experimento de umbrales.
        ¿Qué error es más grave en este contexto y por qué?
  \item ¿Por qué la histéresis evita que la alarma ``parpadee''? Dé un ejemplo numérico.
  \item ¿Por qué no se puede reproducir el sonido de la alarma con \codigo{st.audio} desde el procesador de
        video? ¿Qué cambiaría si la cámara estuviera en otro equipo?
  \item ¿Qué ventajas y desventajas tiene calibrar el umbral por conductor frente a un umbral fijo?
  \item ¿Qué riesgos tiene enviar alertas automáticas? Proponga dos medidas para reducir las falsas alertas.
  \item Compare EAR y CNN en exactitud, velocidad, necesidad de datos y explicabilidad.
  \item ¿Por qué una expresión facial ``negativa'' no indica necesariamente que el conductor esté en riesgo?
  \item ¿Qué sensores adicionales usaría en un vehículo real para detectar un siniestro?
\end{enumerate}

% ======================================================================
\section{Entregables y evaluación}
% ======================================================================
\textbf{Entregables:}
\begin{enumerate}
  \item Notebooks 03 y 04 ejecutados, con las tablas de esta guía completas.
  \item Capturas de pantalla de la aplicación: semáforo en cada nivel, cuenta regresiva, historial y mensaje
        recibido por el contacto.
  \item Archivo CSV de eventos de una sesión en vivo.
  \item Informe (máximo 6 páginas) con resultados, respuestas a las preguntas de análisis y conclusiones.
\end{enumerate}

\begin{table}[H]
\centering
\caption{Rúbrica de evaluación.}
\begin{tabularx}{\textwidth}{>{\raggedright}p{3.4cm} X c}
\toprule
\textbf{Criterio} & \textbf{Descripción} & \textbf{Peso}\\
\midrule
Funcionamiento & Entorno instalado, aplicación en vivo, notificaciones y CNN funcionando. & 30\,\%\\
Evaluación con métricas & Tablas completas: detección, falsas alarmas, latencia, FPS, F1 y matriz de confusión. & 30\,\%\\
Análisis crítico & Explicación de errores, compromiso de umbrales, limitaciones y respuestas a las preguntas. & 25\,\%\\
Documentación & Informe claro, capturas, código comentado y referencias. & 15\,\%\\
\midrule
\multicolumn{2}{r}{\textbf{Total}} & \textbf{100\,\%}\\
\bottomrule
\end{tabularx}
\end{table}

% ======================================================================
\section{Consideraciones éticas y de seguridad}
% ======================================================================
\begin{itemize}
  \item \textbf{Consentimiento:} toda persona que aparezca en la cámara debe autorizarlo.
  \item \textbf{Privacidad:} el sistema no guarda imágenes, solo métricas. Proteja el acceso a MongoDB, que
        contiene datos de contacto.
  \item \textbf{Alertas:} use contactos de prueba que sepan que recibirán mensajes. Nunca envíe alertas a
        servicios de emergencia reales.
  \item \textbf{Expresiones:} no use las etiquetas de expresión para evaluar, calificar ni tomar decisiones
        sobre personas.
  \item \textbf{Seguridad vial:} no pruebe el sistema mientras conduce. Un prototipo puede fallar.
\end{itemize}

% ======================================================================
\section{Solución de problemas}
% ======================================================================
\begin{table}[H]
\centering
\begin{tabularx}{\textwidth}{>{\raggedright}p{4.6cm} X}
\toprule
\textbf{Síntoma} & \textbf{Solución}\\
\midrule
API no disponible & Inicie uvicorn desde la raíz del proyecto (consola 1).\\
MongoDB no disponible & Inicie el servicio \emph{MongoDB Server} o revise \codigo{MONGO\_URI}.\\
No se pudo iniciar el detector & Ejecute \codigo{python scripts/descargar\_modelos.py}.\\
La cámara no inicia & Permita la cámara en el navegador y cierre Teams, Zoom u otras apps.\\
No suena la alarma & Pruebe \codigo{ALARMA\_BACKEND=winsound} y revise el volumen.\\
Muchas falsas alarmas & Recalibre con buena luz frontal; ajuste \codigo{factor\_ear} o \codigo{microsueno\_s}.\\
Error de autenticación en Gmail & Use una contraseña de aplicación, no la contraseña normal.\\
Telegram: \emph{chat not found} & El contacto debe escribir primero al bot.\\
Video lento & Suba \codigo{EXPRESION\_INTERVALO\_S} o desactive las expresiones.\\
\bottomrule
\end{tabularx}
\end{table}

% ======================================================================
\begin{thebibliography}{9}
\bibitem{soukupova2016} T.~Soukupová y J.~Čech, ``Real-Time Eye Blink Detection using Facial Landmarks,''
  en \emph{21st Computer Vision Winter Workshop}, 2016.
\bibitem{wierwille1994} W.~W.~Wierwille y L.~A.~Ellsworth, ``Evaluation of driver drowsiness by trained
  raters,'' \emph{Accident Analysis \& Prevention}, vol.~26, no.~5, pp.~571--581, 1994.
\bibitem{fusek2018} R.~Fusek, ``Pupil localization using geodesic distance,'' en \emph{Advances in Visual
  Computing (ISVC 2018)}, Springer, 2018. (MRL Eye Dataset.)
\bibitem{dosovitskiy2021} A.~Dosovitskiy \emph{et al.}, ``An Image is Worth 16x16 Words: Transformers for
  Image Recognition at Scale,'' en \emph{ICLR}, 2021.
\bibitem{mediapipe} Google, ``MediaPipe Face Landmarker,''
  \url{https://ai.google.dev/edge/mediapipe/solutions/vision/face_landmarker}.
\end{thebibliography}

% ======================================================================
\appendix
\section{Estructura del proyecto}
\begin{lstlisting}
DriveGuardLab/
|-- core/        config, landmarks, riesgo, pipeline, alarma, faces,
|                ojos_cnn, notificaciones, impacto, simulacion, storage
|-- backend/     main, conductores_api, sesiones_api, alertas_api, schemas
|-- frontend/    app, monitoreo_ui, conductor_ui, video_ui, historial_ui, config_ui
|-- notebooks/   00_Verificacion_Entorno, 03_Somnolencia_EAR_MAR,
|                04_CNN_Ojos_Abiertos_Cerrados
|-- scripts/     descargar_modelos, generar_recursos, datos_demo
|-- tests/       43 pruebas (pytest)
|-- docs/        esta guía
|-- assets/  data/  models/
\end{lstlisting}

\section{Endpoints de la API}
\begin{table}[H]
\centering
\begin{tabularx}{\textwidth}{l l X}
\toprule
\textbf{Método} & \textbf{Ruta} & \textbf{Descripción}\\
\midrule
GET & \codigo{/health} & Estado de la API, MongoDB y canales\\
POST/GET & \codigo{/conductores} & Crear y listar conductores\\
GET/PUT/DELETE & \codigo{/conductores/\{id\}} & Consultar, actualizar o eliminar\\
POST/GET & \codigo{/sesiones} & Crear y listar sesiones\\
PUT & \codigo{/sesiones/\{id\}/cerrar} & Cerrar sesión y calcular resumen\\
POST & \codigo{/eventos/lote} & Guardar eventos en lote\\
GET & \codigo{/estadisticas/\{sesion\_id\}} & Resumen de una sesión\\
POST & \codigo{/alertas/emergencia} & Protocolo de emergencia\\
POST & \codigo{/alertas/prueba} & Mensaje de prueba\\
GET & \codigo{/alertas/telegram/chats} & Obtener chat\_id de Telegram\\
\bottomrule
\end{tabularx}
\end{table}

\end{document}
